\documentclass[11pt, a4paper, logo, copyright]{foundationagents}

\pdfinfoomitdate 1
\renewcommand{\today}{}

\usepackage[numbers,sort&compress]{natbib}
\hypersetup{
  colorlinks,
  citecolor=blue,
  linkcolor=red,
  urlcolor=blue}

\usepackage{graphicx}
\usepackage{subcaption}
\usepackage{booktabs}
\usepackage{amsmath}
\usepackage{amssymb}
\usepackage{mathtools}
\usepackage{amsthm}
\usepackage{multirow}
\usepackage{xspace}
\usepackage{algorithm}
\usepackage{algpseudocode}
\usepackage[capitalize,noabbrev]{cleveref}

\newcommand{\method}{\textsc{FoundationAgent}\xspace}

\theoremstyle{plain}
\newtheorem{theorem}{Theorem}[section]
\newtheorem{lemma}[theorem]{Lemma}
\theoremstyle{definition}
\newtheorem{definition}[theorem]{Definition}

\title{A Foundation Agents Paper Template}

\author[1,2]{First Author}
\author[2]{Second Author}
\author[1$\dagger$]{Corresponding Author}

\affil[1]{The Hong Kong University of Science and Technology (Guangzhou)}
\affil[2]{Foundation Agents}

\correspondingauthor{Corresponding Author (\texttt{author@example.com})}
\keywords{Agent learning; evaluation; scalable environments}

\begin{abstract}
This repository provides a compact LaTeX template for Foundation Agents
community papers. It keeps the paper body quiet and readable, while using a
Foundation Agents wordmark in the first-page header and a Google-style technical
report layout. Replace this abstract with the short summary of your work.
\end{abstract}

\begin{document}
\maketitle

\section{Introduction}

This is a minimal paper skeleton using the \texttt{foundationagents} class.
It supports regular sections, citations~\citep{vaswani2017attention}, figures,
tables, theorem environments, algorithms, and appendices.

\section{Method}

The template is intentionally light: keep paper-specific macros in
\texttt{main.tex}, and keep shared layout behavior in \texttt{foundationagents.cls}.
For example, \method can be introduced as a project-local command.

\begin{definition}[Reusable Template]
A reusable template should make the common path easy while leaving paper-specific
styling and notation in the paper source.
\end{definition}

\begin{algorithm}[t]
\caption{Example Procedure}
\label{alg:example}
\begin{algorithmic}[1]
  \State Initialize a candidate solution.
  \For{$t = 1$ to $T$}
    \State Evaluate the current candidate.
    \State Update the candidate using feedback.
  \EndFor
  \State \Return the best candidate.
\end{algorithmic}
\end{algorithm}

\section{Experiments}

\begin{table}[h]
\caption{A compact example table.}
\label{tab:example}
\centering
\begin{tabular}{lcc}
\toprule
Method & Score & Cost \\
\midrule
Baseline & 42.0 & 1.0x \\
\method & 58.5 & 1.2x \\
\bottomrule
\end{tabular}
\end{table}

\section{Conclusion}

Replace this sample with your paper content. The class will place the
Foundation Agents mark in the first-page header when the \texttt{logo} option is
enabled.

\bibliographystyle{plainnat}
\bibliography{references}

\newpage
\appendix
\section{Appendix}

Additional details can go here.

\end{document}
